\documentclass[11pt]{article}

\usepackage[T1]{fontenc}
\usepackage[utf8]{inputenc}
\usepackage{lmodern}
\usepackage[margin=1in]{geometry}
\usepackage{amsmath,amssymb}
\usepackage{booktabs}
\usepackage{microtype}
\usepackage[hidelinks]{hyperref}
\usepackage{enumitem}
\setlist[itemize]{itemsep=2pt,topsep=3pt,parsep=0pt}
\setlist[enumerate]{itemsep=2pt,topsep=3pt,parsep=0pt}

\newcommand{\R}{R}
\newcommand{\K}{K}
\newcommand{\omegared}{\omega_{\mathrm{red}}}
\newcommand{\omegablue}{\omega_{\mathrm{blue}}}

\title{$R(8,10)\ge 346$: a colouring of $K_{345}$ with no red $K_8$ and no blue $K_{10}$}
\author{%
  \normalsize David Lai\\
  \normalsize \texttt{davidlai911@gmail.com}
}
\date{3 November 2026}

\begin{document}
\maketitle

\begin{abstract}
We describe an edge colouring of the complete graph $K_{345}$ containing no red $K_8$ and no blue
$K_{10}$, which proves $\R(8,10)\ge 346$. The previous bound, $\R(8,10)\ge 345$, is due to Cremer
\cite{cremer}, improving the value $343$ recorded in revision~18 of Radziszowski's dynamic survey
\cite{ds1}. Our colouring is obtained by the same device, applied once more: add one vertex to
Cremer's colouring of $K_{344}$ and repair the monochromatic cliques that appear. The added vertex
creates $29\,814$ blue $K_{10}$'s, and the repair reduces that count to zero. We report the colouring,
four independent verifications of it, and the search trajectory, including the one element of the
search that appeared to be decisive: restarting from the best colouring found so far with fresh
random seeds, round after round.
\end{abstract}

\section{Introduction}

Write $\R(r,s)$ for the least $n$ such that every red/blue colouring of the edges of $K_n$ contains a
red $K_r$ or a blue $K_s$. An \emph{$(r,s)$-colouring} of $K_n$ is one containing neither, and shows
$\R(r,s)\ge n+1$. Table~2.1 of Radziszowski's dynamic survey \cite{ds1}, in its current revision,
lists
\[
  343 \;\le\; \R(8,10) \;\le\; 4402 .
\]
The lower bound $343$ is due to Kuznetsov \cite{kuznetsov}, whose witness is a $(8,10)$-colouring of
$K_{342}$ that is circulant: $\{i,j\}$ is red exactly when $\min(|i-j|,342-|i-j|)$ lies in a fixed set
of $76$ distances. Machine search has supplied a number of nearby entries in the same table; Nagda,
Raghavan and Thakurta \cite{nagda}, for instance, improved nine bounds of the form $R(3,s)$ and
$R(4,s)$. Cremer \cite{cremer} observed that adding a single vertex to that colouring, joined
by a blue edge, produces $10\,146$ blue $K_{10}$'s, and that local search can remove all of them; the
resulting colouring of $K_{343}$ proves $\R(8,10)\ge 344$. Adding one further vertex and repeating
the process gives a colouring of $K_{344}$ and the bound $\R(8,10)\ge 345$ \cite{cremer}.

This note records the next step. Starting from Cremer's colouring of $K_{344}$, adding one vertex
creates $29\,814$ blue $K_{10}$'s --- the exact minimum over all $2\cdot 344$ choices of which vertex
to copy and which colour to give the joining edge --- and exact-conflict local search reduces that
number to zero. The result is a $(8,10)$-colouring of $K_{345}$:

\begin{center}
\fbox{\parbox{0.92\textwidth}{
\textbf{Theorem.} $\R(8,10)\ge 346$.
}}
\end{center}

Section~\ref{sec:colouring} states the colouring and its basic invariants;
Section~\ref{sec:verification} records four independent verifications;
Section~\ref{sec:method} describes how the colouring was found and what in the search mattered;
Section~\ref{sec:ai} describes the AI system that ran the search and what it contributed;
Section~\ref{sec:not} says what the computation does not establish; and
Section~\ref{sec:data} lists the artefacts needed to reproduce the result; and
Appendix~\ref{app:distance} records an exact bounded-distance screen of the same one-vertex step on
six smaller cells, which shows how far a repair must travel before it can succeed at all.

\paragraph{Provenance.}
The search code was written and run by an AI agent under the author's direction, following the same
arrangement described in Section~5 of \cite{cremer}. The screening stage was a multi-agent run; the
final descent was not. Section~\ref{sec:ai} describes both, and what each contributed. Every number
in this note was regenerated and checked independently of the search process before being recorded
here; Section~\ref{sec:verification} specifies the checkers, and Section~\ref{sec:data} lists the
artefacts needed to reproduce or refute the result.

\section{The colouring}\label{sec:colouring}

The colouring is stored as a $345\times 345$ matrix $M$ with entries in $\{0,1\}$, symmetric, with
zero diagonal, where $M_{ij}=1$ means that $\{i,j\}$ is red. Its invariants are as follows.

\begin{center}
\begin{tabular}{lr}
\toprule
vertices & $345$\\
red edges & $26\,045$\\
blue edges & $33\,295$\\
$\omegared$ (largest red clique) & $7$\\
$\omegablue$ (largest blue clique) & $9$\\
red degrees & $146$ to $160$, mean $150.99$\\
distinct red degrees & $15$\\
\bottomrule
\end{tabular}
\end{center}

Since $\omegared=7<8$ and $\omegablue=9<10$, the colouring contains no red $K_8$ and no blue
$K_{10}$, so $\R(8,10)\ge 346$.

As in \cite{cremer} we note that the red graph is not circulant: its red degrees take $15$ distinct
values, so no relabelling makes it invariant under a cyclic shift. Nor is it a near-copy of any
colourings the search passed through: it differs from the $6$-conflict state it was built from in
$1\,032$ edges and from a $1$-conflict state found alongside it in $1\,226$ (Section~\ref{sec:method}).

\section{Verification}\label{sec:verification}

The colouring was verified by four programs, written independently of one another and of the search
code.

\begin{center}
\small
\begin{tabular}{@{}lll@{}}
\toprule
checker & method & result\\
\midrule
\texttt{count} & exhaustive enumeration of red $K_8$, blue $K_{10}$ & $0$, $0$\\
\texttt{verify} & maximum-clique branch and bound & $\omegared=7$, $\omegablue=9$\\
\texttt{verify2} & target-size DFS, independent code & no red $K_8$, no blue $K_{10}$\\
\texttt{enum} & independent clique enumerator & $0$ conflicts\\
\bottomrule
\end{tabular}
\end{center}

The two count-based checkers agree with the two structural ones: a colouring with $\omegared=7$ and
$\omegablue=9$ contains no forbidden clique, and the counts confirm the same directly. To indicate the
cost of an independent re-check, \texttt{verify2} searched $10\,471\,027$ nodes for a red $K_8$ and
$195\,066\,549$ for a blue $K_{10}$ on a single core; the blue search is two orders of magnitude
larger, as expected when the forbidden blue clique is two vertices bigger.

The files were also decoded from their graph6 representation \cite{networkx} and compared edge by
edge with the matrix form, so that the two deposited encodings are known to carry the same graph.

\section{How the colouring was found}\label{sec:method}

\paragraph{The setting.}
A \emph{conflict} is a monochromatic clique of a forbidden size: a red $K_8$ or a blue $K_{10}$.
The engine used here was written during the screening stage described in \S\ref{sec:ai}.
The search keeps the complete set of conflicts, together with the number of conflicts lying on each
edge, so that the effect of flipping any edge is computed exactly rather than sampled. One iteration
chooses a conflict at random with probability proportional to $w\,|K|$ over the two colours, samples
up to $32$ of its edges, and flips the edge of least weighted exact change $\delta$, accepting it if
$\delta\le 0$ or with Metropolis probability $\exp(-\delta/T)$, where $T=0.25+\min(2.5,t/100)$ and
$t$ counts iterations since the last improvement; the flipped edge is made tabu for $5$ to $19$
iterations. The weights are $w_{\mathrm{red}}=\binom{10}{2}/\binom{8}{2}=45/28$ and
$w_{\mathrm{blue}}=1$. When no improvement has occurred for $1\,000$ iterations, the search restores
the best colouring found so far, performs eight conflict-directed random flips, clears the tabu list
and re-randomises the red weight; in addition, a \emph{patch} move occasionally randomises the colours
of pairs inside a set of up to eight vertices chosen around a vertex with probability proportional to
its conflict count, and then repairs freely.

\paragraph{The first step, exactly.}
Copying a vertex $v$ and joining the copy $v'$ to $v$ by an edge of colour $c$ creates a conflict
$\{v,v'\}\cup C$ for each $c$-clique $C$ of size $6$ (for red) or $8$ (for blue) inside $N_c(v)$, and
nothing else. One clique pass per vertex and colour therefore gives all $2n$ twin costs exactly. For
Cremer's $K_{344}$ colouring the minimum over all $688$ choices is $29\,814$, attained by the blue
twin of vertex $335$; the cheapest red twin costs $41\,479$. The spread between twins is large, so
this choice is worth making exactly.

\paragraph{The descent.}
From $29\,814$ conflicts the search descended as follows. Each row is a round of one or more
single-threaded segments; within a round the best colouring of the previous round is the input, and
the random seeds are new.

\begin{center}
\begin{tabular}{@{}rlll@{}}
\toprule
conflicts & round & lanes & wall clock\\
\midrule
$29\,814 \to 97$ & pilot, 13 segments, 11 seeds & 1 & $\approx 102$ CPU-minutes\\
$97 \to 96$ & single 90-second segment & 1 & $90$ s\\
$96 \to 41$ & round 1 & 4 & $\approx 2$ h\\
$41 \to 23$ & round 2 & 4 & $\approx 2$ h\\
$23 \to 6$ & round 3 & 4 & $\approx 2$ h\\
$6 \to 0$ & round 4 & 4 & $\approx 2$ h\\
\bottomrule
\end{tabular}
\end{center}

The colouring that achieves zero conflicts comes from the segment that started at the $6$-conflict
state and reached zero at iteration $159\,966$ after $2\,770$ seconds, its conflict count falling as
$6\to5\to4\to3\to2\to1\to0$. The colours moved a long way in the process: the final colouring
differs from the $6$-conflict state it started from in $1\,032$ edges, and from a $1$-conflict
colouring found in a parallel lane in $1\,226$ edges. Two colourings that are each within one
monochromatic clique of being valid can thus be separated by more than a thousand edge changes, so
the successful colouring is not a small repair of anything the search had already seen. For
comparison, the same engine had earlier reduced an $11$-conflict state to $2$ and then made no further
progress for two thousand seconds.

\paragraph{What mattered.}
Two things, in our assessment.

First, \emph{restarts from the best state with fresh seeds}. A single long segment is misleading here.
In one run the improvement rate collapsed to $0.003$ conflicts per $100$ seconds at $97$ conflicts, a
profile that reads as convergence to a floor. It was not: restarting from that same colouring with new
seeds reached $41$ conflicts within two hours, and three further rounds reached zero. Every round of
fresh seeds found a deeper minimum than the last. Any assessment of this kind of search should be made
per round, not per segment.

Second, the \emph{last conflicts are not the cheapest}. One lane of the final round reached a
$2$-conflict state consisting of a single red $K_8$ and a single blue $K_{10}$; an exhaustive check of
all $73$ edges of those two cliques, together with the intersection condition for a second flip,
showed that no single flip and no pair of flips removes both, so a repair had to change at least three
coordinated edges. Breaking a red conflict turns those edges blue and so creates blue conflicts, and
conversely, which is why the two ends pull against each other. The descent from $29\,814$ to $6$ was
therefore not the difficult part: it was the last few conflicts that required the search to leave the
neighbourhood it had been exploring, and the successful segment spent $2\,770$ seconds and
$159\,966$ iterations on its final six.

Conversely, we found no use for a donor-crossover operator at this scale. Two segments running a
section-crossover variant, given the same budget from the same states, failed to improve on the plain
patch move at $96$ conflicts and at $6$ conflicts, while the plain lanes improved at both. This is
consistent with \cite{cremer}, where crossover is credited with closing a state at $6$--$7$ conflicts
after the repair stage had plateaued, but we did not reproduce that advantage here.

\section{On the use of AI}\label{sec:ai}

\paragraph{The system.}
The computational work was carried out by \emph{Open Math Model}, running DeepSeek V4.1 Flash --- a
model supplied by DeepSeek --- at maximum reasoning effort, operating an ordinary Linux virtual
machine with four cores. The
agent wrote all programs, chose the search strategy, ran the experiments and did the bookkeeping; the author set the goal, fixed the compute budget and decided
which cells to pursue. This is the same division of labour reported in \cite{cremer}.

\paragraph{The screening swarm.}
The first stage ran as an \emph{Ultramath} swarm: eight identical agents (labelled \texttt{m1} to
\texttt{m8}) attacking one shared problem concurrently. Each was required to register a distinct line
of attack before starting --- a duplicate registration is refused --- and the agents exchanged
findings on a common board rather than working in isolation. The problem given to the swarm was to
apply the add-a-vertex-then-repair method of \cite{cremer} to as many cells of Table~2.1 as possible
and to determine, by measurement rather than assertion, which cells it can move. The run lasted four
hours with a budget of one hundred model steps per member, on the same four cores used later for the
present computation.

\begin{center}
\small
\begin{tabular}{@{}ll@{}}
\toprule
system & Open Math Model\\
model & DeepSeek V4.1 Flash\\
supplier & DeepSeek\\
reasoning effort & maximum\\
swarm size & $8$ agents, one shared problem, $10$ distinct directions opened\\
duration & $4$ hours on $4$ cores\\
outcome & $8$ submissions, all accepted after independent checking\\
\bottomrule
\end{tabular}
\end{center}

\paragraph{What the swarm contributed.}
Four things, all of which are used in this paper.

\begin{enumerate}
\item An \emph{exact twin-insertion table} for nine cells --- for every vertex $v$ and both colours
      of the joining edge, the number of conflicts created by copying $v$ --- covering the record
      colourings then available for those cells, several of which are due to Exoo and Tatarevic
      \cite{exoo-tatarevic}. This is the quantity that
      determines the starting point of \S\ref{sec:method}, and its minimum over the $688$ choices for
      the $(8,10)$ cell is $29\,814$.
\item \emph{Reusable engines.} The exact-conflict repair program used for the descent in
      \S\ref{sec:method} was written during the swarm, as was a second engine implementing the
      preparation, risk-chain and crossover stages of \cite{cremer}, a depth-bounded exhaustive
      repair checker, and a twin-scanning utility.
\item A \emph{bounded-distance screen}: a CEGAR-SAT encoding that decides, for increasing $k$,
      whether any valid colouring of the target order lies within Hamming distance $k$ of a given
      $+1$ state. Applied to six cells, it proved that no valid colouring lies within $9$ to $17$
      edge changes of the corresponding state --- a statement no heuristic can make --- and thereby
      established that those cells cannot be moved from those starting colourings by this method at
      all, however long it runs. The screen is reproduced in Appendix~\ref{app:distance}.
\item A \emph{negative result about search heuristics}: the conflict count of the $+1$ state does not
      rank how far a cell is from a solution. One cell with $72$ conflicts after the $+1$ step needed
      only $9$ edge changes, while another with $6$ needed $13$. Every agent in the swarm, and the
      author's own initial reading, ranked the cells by conflict count; that ranking was wrong, and
      the $(8,10)$ cell --- which looked among the least promising, with $29\,814$ initial
      conflicts --- is the one that moved.
\end{enumerate}

\paragraph{The final campaign.}
The swarm finished with no cell moved. The descent reported in \S\ref{sec:method} was run afterwards
by the agent alone, using the engine from item~2 above, in six rounds of four parallel single-threaded
lanes over approximately eight hours of wall-clock time. The AI's role in the two stages differed in
kind: the swarm was useful for breadth and for producing tools and negative results, while the
positive result came from a long, narrow, single-purpose campaign that no member had the budget to
run and that the swarm's own screening had not identified as the most likely to succeed.

\paragraph{Verification by the agent, not by the model.}
Every numeric claim in this paper was recomputed after the fact, outside the search process and
largely outside the code that produced it, before being recorded: the witness by four independent
checkers, the twin tables by three separate implementations, the distance bounds against positive
controls at a known distance, and the search trajectories against exact conflict counts. Two errors
in the intermediate records were caught this way. A depth-bounded exhaustive repair checker reported a
repair from a $2$-conflict state at $n=345$ and wrote back an unchanged matrix; an independent
enumeration of all $73$ candidate first flips showed that no $1$- or $2$-flip repair existed there,
and that checker's large-$n$ output is not used anywhere in this paper (it is reliable against
controls at $n=40$ and $n=58$). Separately, an earlier reading of the search described a descent that
had stalled at $97$ conflicts as approaching a floor; it was seed exhaustion, and the descent resumed
as soon as the seeds were renewed, as described in \S\ref{sec:method}.

\section{What the computation does not show}\label{sec:not}

\begin{itemize}
\item The upper bound is unchanged; $\R(8,10)\le 4402$ \cite{ds1} is far from $346$. Nothing here
      bears on the value of $\R(8,10)$, only on its lower end.
\item No statement is made about the existence of $(8,10)$-colourings of $K_{346}$. A failed search
      for one would not be evidence of its non-existence, and we did not run that search.
\item The colouring is not presented as canonical or as one of a known family. It was found by
      stochastic search; whether it is isolated among $345$-vertex colourings is open.
\item The bounded-distance screen of Appendix~\ref{app:distance} produced no bound at any of the
      cells it covers, and it did not predict the present result: no such screen finishes at
      $n=345$, so the $(8,10)$ cell was outside what it could decide in either direction.
\end{itemize}

\section{Reproduction}\label{sec:data}

\paragraph{Artefacts.}
The colouring is deposited in two encodings:
\begin{center}
\small
\begin{tabular}{@{}ll@{}}
\toprule
matrix, $345$ rows of \texttt{0}/\texttt{1} ($1$ = red) & \texttt{r810-n345-k8-free-k10-free.txt}\\
graph6 of the red graph & \texttt{r810-n345-k8-free-k10-red.g6}\\
\midrule
SHA-256 of the matrix & \texttt{0af4590108642651e3999994b379235}\\
 & \texttt{f8b126bb9eb1fdfd36fb450ef47935b12}\\
SHA-256 of the graph6 file & \texttt{4135f58f2cb289bfa0ab02ca024e1a2}\\
 & \texttt{d7e5f695aa505e96fa88b4bde54158d62}\\
\bottomrule
\end{tabular}
\end{center}

Splitting each hash across two lines is a typesetting device only; the two halves are one string.

\paragraph{Search.}
The search is single-threaded C. The run that produced the colouring was invoked as
\begin{center}
\texttt{engine start.txt 8 10 0 run 20000000 32 222 3600 1000 out 1.607 64}
\end{center}
from a $345$-vertex input with $6$ conflicts, where the arguments are, in order, the input matrix, the
forbidden red and blue clique sizes, a colour-convention flag, the mode, an iteration cap, the
candidate sample size, the random seed, a CPU-second budget, the stall threshold, the output prefix,
the red weight and a cap on the number of conflicts a flip may create. Inputs are matrices with a
zero diagonal and $1$ denoting a red edge; the colour convention must be stated explicitly, since
published witnesses differ in it, and a matrix inverted by hand to change convention acquires a
non-zero diagonal and will be mis-counted by any clique enumerator that assumes otherwise.

\appendix

\section{How far a repair must travel: an exact bounded-distance screen}\label{app:distance}

The method studied in this paper is a repair: from a colouring that is not yet valid, change edges
until no forbidden clique remains. Whether that is possible at all is a question about distance, and
inside a Hamming ball the question is decidable exactly. This appendix records the answers for six
smaller cells, obtained during the screening stage described in \S\ref{sec:ai}. They are the reason
the present cell was the one worth pushing, and they delimit the method in a way that no search can.

\subsection*{The question}

Let $B$ be a colouring of $K_{n_0+1}$ obtained by adding one vertex to a valid $(r,s)$-colouring of
$K_{n_0}$ --- a twin, in the sense of \S\ref{sec:method}. For $k=0,1,2,\dots$ ask:

\begin{center}
\emph{is there an $(r,s)$-colouring of $K_{n_0+1}$ at Hamming distance at most $k$ from $B$?}
\end{center}

An unsatisfiable instance at bound $k$ is a proof that \emph{no} sequence of at most $k$ edge changes
repairs $B$, whatever heuristic is applied. It is therefore a statement about the cell, not about a
program.

\subsection*{The encoding}

One Boolean variable per edge, with $1$ meaning red. The constraints are that no red $K_r$ and no blue
$K_s$ exist. The red side is usually small enough to write out in full; the blue side is not --- for
$n=40$ there are $\binom{40}{10}=8.5\times10^{8}$ ten-sets --- so the clique constraints are added
lazily: solve, enumerate every violated clique exactly, add the clause forbidding that clique, and
re-solve. The distance bound is imposed by a totalizer over flip variables
$d_e \leftrightarrow (x_e \neq B_e)$, asserting $\sum_e d_e \le k$.

Two solvers were used, CaDiCaL~1.5.3 and MiniSat~2.2, and they agree on every bound that both
completed.

\subsection*{A positive control is mandatory}

An encoding that is too strong returns ``unsatisfiable'' for every input, which looks identical to a
correct proof. The screen was therefore validated against a case whose answer is known: take a valid
$(4,6)$-colouring of $K_{35}$, flip $5$ of its edges at random, and require the encoder to return
satisfiable at $k=5$ with a witness that re-verifies to zero conflicts, and unsatisfiable at $k=4$.
Both behaviours were observed, and the returned witness was checked by an independent enumerator.
Without such a control a distance bound is worthless.

\subsection*{Results}

The column $k^\ast$ is the largest $k$ for which unsatisfiability was proved, so a repair from that
state needs at least $k^\ast+1$ edge changes.

\begin{center}
\small
\begin{tabular}{@{}llrrl@{}}
\toprule
cell & best twin & conflicts after $+1$ & $k^\ast$ & a repair needs\\
\midrule
$(5,6)$ & $v_{56}$, red & $136$ & $16$ & $\ge 17$ edge changes\\
$(4,6)$ & $v_{31}$, red & $24$ & $16$ & $\ge 17$\\
$(3,10)$ & $v_{35}$, red & $6$ & $12$ & $\ge 13$\\
$(4,7)$ & $v_{15}$, red & $40$ & $12$ & $\ge 13$\\
$(3,11)$ & $v_{20}$, red & $8$ & $10$ & $\ge 11$\\
$(4,9)$ & $v_{0}$, red & $72$ & $8$ & $\ge 9$\\
\midrule
$(6,8)$, $(6,9)$, $(8,10)$ & --- & $338$, $2\,830$, $29\,814$ & none & undecided\\
\bottomrule
\end{tabular}
\end{center}

For $(3,10)$ a second, better starting state was also screened: a $40$-vertex colouring with four red
triangles and no blue $K_{10}$, obtained by choosing the new vertex's red neighbourhood rather than
copying a vertex. No valid colouring lies within $10$ edge changes of it.

\subsection*{Two conclusions}

\emph{The method cannot move those six cells from those bases.} For $(4,6)$ this is decisive: every
engine in the study, including one running the complete pipeline of \cite{cremer} with preparation,
risk chains and mutation, stalls at $12$ conflicts, and the nearest valid $36$-vertex colouring is at
least $17$ edge changes away. No amount of tuning, restarting or extra seeds closes a gap of $17$.

\emph{Conflict count does not rank repairability.} The $(4,9)$ start has $72$ conflicts and needs only
$9$ edge changes; the $(3,10)$ start has $6$ --- the smallest in the table, and the natural first
choice --- and needs $13$. Ranking these cells by how many conflicts the vertex addition creates, which
is the obvious heuristic and the one adopted at the outset of this work, orders them wrongly. The
$(8,10)$ cell, which looked the least promising of all with $29\,814$ initial conflicts, is the one
that moved.

\subsection*{Limits}

The bounds are for the specific best-twin extension of each base colouring; they are not statements
about the cells, and a colouring of the target order might exist far from those particular states. For
$(6,8)$, $(6,9)$ and $(8,10)$ no bound was completed --- no rung finishes at those sizes --- so for
those cells nothing is claimed in either direction, including for the result of this paper.


\begin{thebibliography}{9}\small

\bibitem{cremer}
F.~Cremer.
\newblock Lower bounds for Ramsey numbers: $\R(6,8)\ge 135$ and $\R(8,10)\ge 345$.
\newblock Preprint, arXiv:2610.12122 [math.CO], 8 October 2026.
\newblock Colourings at \url{https://github.com/fritzcremer/ramsey-lower-bounds}, release v1.1.0.

\bibitem{ds1}
S.~P. Radziszowski.
\newblock Small Ramsey numbers.
\newblock \emph{Electronic Journal of Combinatorics}, Dynamic Survey DS1, revision 18
  (24 April 2026), 149 pp.

\bibitem{kuznetsov}
E.~Kuznetsov.
\newblock Computational lower limits on small Ramsey numbers.
\newblock Preprint, arXiv:1505.07186v5 [cs.DM], 2016.

\bibitem{exoo-tatarevic}
G.~Exoo and M.~Tatarevic.
\newblock New lower bounds for 28 classical Ramsey numbers.
\newblock \emph{Electronic Journal of Combinatorics}, 22(3):Paper P3.11, 12 pp., 2015.
  arXiv:1504.02403.

\bibitem{nagda}
A.~Nagda, P.~Raghavan, and A.~Thakurta.
\newblock Reinforced generation of combinatorial structures: Ramsey numbers.
\newblock Preprint, arXiv:2603.09172v5 [math.CO], 2026.

\bibitem{networkx}
A.~A. Hagberg, D.~A. Schult, and P.~J. Swart.
\newblock Exploring network structure, dynamics, and function using NetworkX.
\newblock In \emph{Proceedings of the 7th Python in Science Conference (SciPy 2008)},
  pages 11--15, 2008.

\end{thebibliography}

\end{document}
